% `tex.web` §320 `macro_call` and §401 `@<Show the text of the macro being
% expanded@>`: what `\tracingmacros` says.
%
%     <blank line>
%     \b #1->[#1]
%     #1<-x
%
% §401's `print_ln`, then §262 `print_cs` — which gives a multi-letter name, and
% a single-letter one whose character is a letter, a trailing space — then
% `token_show`, which is the parameter text, `->`, and the body. Each argument
% follows on its own line as §400 captures it, delimiters and all.
%
% `\meaning` already builds `macro:#1->[#1]` (§296 is `print_cmd_chr` and then
% the same `token_show`), so the trace is that string from after the `macro:`.
%
% Two things this took a probe to get right. An argument lives in the parameter
% stack, not the arena where bodies do — reading the body's range gave every
% argument the text of whatever had been defined last. And §401 opens a line
% only if one is in progress, so the *first* trace of a run has no blank ahead
% of it while every later one does.
%
% Level 2 adds §323 `begin_token_list`: a token list that is not a macro
% announces itself the same way, `\everypar->P`, with no arguments and no blank
% line — §323 opens with `print_nl("")` where §401 opens with `print_ln`.
\catcode`\{=1 \catcode`\}=2 \catcode`\#=6
\font\rip=cmr10 \rip
\hsize=100pt \vsize=50pt \parindent=0pt \hbadness=10000
\def\a{A}
\def\b#1{[#1]}
\def\c#1#2{(#1|#2)}
\def\d#1,#2.{<#1;#2>}
\long\def\e#1{L#1}
\def\f{}

% --- level 1: macro calls and their arguments ------------------------------
\tracingmacros=1
\a
\b{x}
\c{p}{q}
\d y,z.
\e{w}
\f
% an argument that is itself a group, and one that is a control sequence
\b{{g}}
\b\relax
\tracingmacros=0

% --- level 2: token lists that are not macros ------------------------------
\everypar{P}
\output={\shipout\box255}
\tracingmacros=2
\a
X\par
\tracingmacros=0
\end
